\chapter{Contexte Général et Cadre du Projet}

Ce premier chapitre pose les bases de notre environnement de travail. Il présente en premier lieu le cadre institutionnel dans lequel s'inscrit ce projet, avant de détailler l'équipe et le contexte du projet. Enfin, il aborde la méthodologie de gestion de projet adoptée pour assurer le bon déroulement de nos travaux.

\section{Organisme d'accueil}
Ce projet s'inscrit dans le cadre du Projet de Fin d'Année (PFA) de la quatrième année du cursus d'ingénieur à l'\textbf{École Nationale d'Intelligence Artificielle de Berkane (ENIAD)}.

L'ENIAD est une institution de premier plan spécialisée dans la formation d'ingénieurs experts en intelligence artificielle, en science des données et en technologies de l'information. En favorisant l'innovation, la recherche et l'application pratique des nouvelles technologies, l'établissement nous offre un environnement académique propice au développement de solutions technologiques répondant à des problématiques réelles, telles que l'intégration de l'IA dans les domaines de la santé et du bien-être.

\section{Cadre du projet}
Ce projet collaboratif a été réalisé par un trinôme d'élèves ingénieurs en quatrième année : \textbf{BOUDA Mohamed}, \textbf{BOUISIFI Adam} et \textbf{BOURASS Weame}. L'ensemble de notre travail a été supervisé et encadré par notre professeur, \textbf{Pr. MELLAH Youssef}.

L'initiative de ce PFA naît d'une volonté d'appliquer nos connaissances en intelligence artificielle et en développement logiciel sur un cas d'usage concret : le fitness intelligent. Au lieu d'aborder ces deux domaines séparément, notre équipe s'est fixé pour défi de fusionner le coaching sportif et le machine learning dans une même application mobile, offrant ainsi une alternative numérisée, intelligente et accessible par rapport aux salles de sport classiques.

\section{Méthodologie de conduite de projet}
Afin de garantir le succès de ce projet de fin d'année et de respecter nos contraintes de temps, nous avons opté pour la méthode \textbf{Agile Scrum}. Cette approche itérative et incrémentale est particulièrement bien adaptée au développement logiciel car elle favorise la flexibilité, la communication continue au sein de l'équipe et l'adaptation rapide aux changements.

\subsection{L'approche Agile Scrum}
Notre cycle de développement a été divisé en plusieurs \textit{Sprints} de courtes durées (généralement d'une à deux semaines). Chaque Sprint nous a permis de livrer des fonctionnalités opérationnelles de l'application (par exemple : le maquettage, l'intégration du modèle d'IA, l'interface utilisateur). Nous organisions des réunions régulières pour faire le point sur l'avancement, identifier les éventuels blocages et redéfinir nos priorités.

\subsection{Outils de planification et de suivi}
Pour mettre en pratique cette méthodologie, nous nous sommes appuyés sur des outils modernes de gestion de projet :
\begin{itemize}
    \item \textbf{Trello} : Cet outil nous a permis de créer un tableau de bord visuel basé sur le système Kanban. Chaque tâche y était décrite, assignée à un membre de l'équipe et dotée d'une date limite, ce qui a grandement facilité notre coordination.
    \item \textbf{GitHub} : En complément de Trello, GitHub a été essentiel pour le versionnage de notre code, le travail collaboratif simultané sur l'application et la révision mutuelle de nos développements.
\end{itemize}